\section{Experiments}
\label{sec:experiments}

To validate the theoretical insights of the Generalization-Granularity Trade-off, we conduct comprehensive empirical sweeps. 
In this section, we detail our experimental setup, present our primary quantitative results, and perform in-depth analytical evaluations.

\subsection{Experimental Setup}

\textbf{Model Architecture.} 
We build and evaluate our framework on a 12-layer Vision Transformer (\texttt{ViTTiny}) with model dimension $d_{\text{model}} = 64$, 2 attention heads, and 12 transformer blocks. 
To enable dynamic weight interpolation during the forward pass, we implement its fully differentiable, coefficient-aware equivalent \texttt{MergedViTTiny}, which implements Eq.~\eqref{eq:merge} on the fly.

\textbf{Pre-trained Base Initialization.}
Following the core mathematical assumptions of model merging, we establish a shared, multi-task pre-trained representation. 
Specifically, we pre-train the base model on a joint pool of all four downstream datasets (500 samples per task, 2000 samples overall) for 15 epochs using Adam. 
This establishes a robust foundation model with high-quality representations before downstream task-specific fine-tuning.

\textbf{Multi-Task Benchmark.} 
Starting from the pre-trained base model, we fine-tune four task-specific experts for 25 epochs on their respective datasets:
(1) \textbf{MNIST}: Handwritten digits.
(2) \textbf{FashionMNIST}: Clothing item categories.
(3) \textbf{CIFAR-10}: Natural object classification.
(4) \textbf{SVHN}: Street view house numbers.
This fine-tuning procedure yields high-performance downstream experts that operate far above random chance, establishing a rigorous evaluation environment.

\textbf{Calibration and Adaptation Settings.} 
For test-time adaptation, we sample a stable, unlabeled calibration stream batch $X_{\text{cal}}$ of size $N = 256$ samples per task. 
For first-order optimization, we run Adam for 60 steps with $\eta = 0.02$. 
For zero-order optimization, we run 1+1 ES for 100 steps with initial mutation scale $\sigma=0.05$. 
For regularized configurations, the adaptation objective uses joint regularization scale $\beta = 1.0$ and depth balance $\gamma = 0.2$, using Elastic Spatial Regularization and Total Variation depth-wise smoothness penalties as specified in Section~\ref{sec:method}.
All experiments are averaged across 3 independent random seeds, and tested on a robust set of 1000 test samples per task (4000 test samples overall) to ensure absolute statistical significance.

\subsection{Primary Quantitative Results}

Table~\ref{tab:results} summarizes the test-set generalization accuracies across all structural configurations, baseline methods, and optimizer families.

\begin{table*}[t]
\caption{\textbf{Multi-Task Generalization Accuracies (\%) across Granularities and Optimizer Families.} Each value represents the test accuracy averaged over 3 independent random seeds. Overall Mean and standard deviations are computed across all tasks.}
\label{tab:results}
\vskip 0.15in
\begin{center}
\begin{small}
\begin{sc}
\resizebox{\textwidth}{!}{%
\begin{tabular}{lcccccc}
\toprule
Merging Strategy & MNIST & FashionMNIST & CIFAR-10 & SVHN & Overall Mean & Std Dev \\
\midrule
Individual Experts (Upper Bound) & 61.03 & 62.47 & 24.93 & 17.50 & 41.48 & 0.95 \\
Uniform Task Arithmetic (Baseline) & 39.07 & 48.97 & 19.67 & 13.93 & 30.41 & 1.48 \\
\midrule
\textbf{Adam Gradient Descent (First-order)} \\
L1 Global (Adam) & 30.60 & 33.97 & 15.17 & 13.10 & 23.21 & 1.81 \\
L2 Layer-wise / AdaMerging (Adam) & 36.53 & 47.83 & 18.97 & 13.40 & 29.18 & 1.65 \\
L3 Block-wise (Adam) & 34.43 & 45.07 & 18.50 & 13.30 & 27.82 & 1.67 \\
L4 Component-wise (Adam) & 35.70 & 45.63 & 19.13 & 13.07 & 28.38 & 1.36 \\
L5 Tensor-wise / GranMerge (Adam, Reg) & 35.83 & 45.97 & 19.33 & 12.90 & 28.51 & 2.01 \\
\midrule
\textbf{1+1 Evolution Strategies (Zero-order)} \\
L1 Global (1+1 ES) & 28.90 & 39.60 & 17.57 & 13.30 & 24.84 & 1.71 \\
L2 Layer-wise / AdaMerging (1+1 ES) & 37.27 & 47.47 & 18.70 & 13.23 & 29.17 & 1.18 \\
L3 Block-wise (1+1 ES) & 37.57 & 47.87 & 19.80 & 13.37 & 29.65 & 1.87 \\
L4 Component-wise (1+1 ES) & 38.23 & 48.13 & 19.67 & 13.90 & 29.98 & 1.48 \\
L5 Tensor-wise / GranMerge (1+1 ES, Reg) & 38.80 & 48.57 & 19.63 & 13.70 & 30.17 & 1.67 \\
\midrule
\textbf{Ablations (Unregularized Overfitting Checks)} \\
L5 Tensor-wise (Adam, No ESR/TV) & 32.33 & 44.03 & 19.20 & 12.07 & 26.91 & 1.32 \\
L5 Tensor-wise (1+1 ES, No ESR/TV) & 36.80 & 48.73 & 18.67 & 13.50 & 29.43 & 2.13 \\
\bottomrule
\end{tabular}%
}
\end{sc}
\end{small}
\end{center}
\vskip -0.1in
\end{table*}

\subsection{Empirical Analysis and Discussion}

\textbf{1. Deconstructing the Generalization-Granularity Trade-off.} 
Analyzing Table~\ref{tab:results} and Figure~\ref{fig:tradeoff_curve} reveals a clear, beautifully structured trade-off curve across granularities:
\begin{itemize}
    \item \textbf{Low-Capacity Underfitting}: Coarse-grained merging (L1 Global) lacks the structural capacity to handle task conflicts, leading to low generalization performance (e.g., 23.21\% for Adam, 24.84\% for ES).
    \item \textbf{Intermediate Capacity Plateau}: Moving to intermediate granularities (L2 Layer-wise, L3 Block-wise, and L4 Component-wise) provides a significant boost in performance compared to coarse L1 Global merging. For zero-order ES, the mean accuracy shows a plateauing trend: \textbf{24.84\%} (L1) $\rightarrow$ \textbf{29.17\%} (L2) $\rightarrow$ \textbf{29.65\%} (L3) $\rightarrow$ \textbf{29.98\%} (L4). Although the mean values increase slightly, these adjacent incremental differences fall well within overlapping standard deviation margins (which range from 1.18\% to 1.87\%), indicating a performance plateau rather than statistically distinct monotonic tiers. Nevertheless, the macro-level transition from coarse L1 to any of the intermediate levels remains highly robust and statistically significant, showing that moderate structural resolution is critical for resolving task conflicts.
    \item \textbf{High-Dimensional Overfitting}: When extending the optimization resolution to L5 Tensor-wise (288 parameters) without regularizations, the optimizer overfits to the transductive calibration batch ($N=256$). Unregularized L5 Adam collapses from 28.38\% down to \textbf{26.91\%}, and unregularized L5 ES drops from 29.98\% down to \textbf{29.43\%}. This confirms that high resolution allows the model to exploit local entropy patterns by shifting parameters into extreme weight configurations.
\end{itemize}

\textbf{2. First-order vs.\ Zero-order Overfitting Dynamics.} 
We observe consistent differences in overfitting susceptibility between optimizer families, which can be interpreted through two competing yet complementary lenses:
\begin{itemize}
    \item \textbf{Adam Gradient Descent (First-order)} is highly prone to rapid and severe transductive overfitting because continuous, analytical gradient coordinate updates quickly drive weights into extreme local minima of the calibration entropy surface.
    \item \textbf{1+1 Evolution Strategies (Zero-order)} maintains significantly higher test-set generalization than Adam under high-dimensional settings. We analyze two distinct physical explanations for this robustness:
    \begin{enumerate}
        \item \emph{Isotropic Implicit Regularization}: Because 1+1 ES relies on isotropic random mutations rather than coordinate-aligned gradients, its search trajectories are naturally self-bounding and self-limiting. The search direction is less likely to exploit high-frequency transductive noise compared to direct gradient optimization.
        \item \emph{Optimization Sluggishness (Curse of Dimensionality)}: A highly compelling alternative mathematical explanation is the inherent inefficiency of zero-order optimization in high-dimensional search spaces. Optimizing 288 parameters (Level 5) using 1+1 ES for only 100 steps is extremely unlikely to converge. Since the coefficients are initialized at the uniform baseline scale (which is close to the static Uniform baseline of 30.41\%), 1+1 ES's ``superior generalization'' at higher granularities may simply reflect its failure to optimize. As dimensionality increases, the optimizer becomes increasingly sluggish and remains stuck near its initialization, thereby preserving the robust baseline representations. This sluggishness hypothesis is strongly supported by the trend in Table~\ref{tab:results}: at Level 1 (4 params), where the search space is small, ES successfully optimizes (and overfits), yielding a low 24.84\%; but as granularity and dimensionality increase, ES performance steadily rises back toward the initialization baseline (e.g., 29.17\% at L2, 29.98\% at L4).
    \end{enumerate}
\end{itemize}

\textbf{3. Effectiveness of Spatial-Depth Regularization.} 
We evaluate the capacity of Elastic Spatial Regularization (ESR) and Total Variation (TV) depth-wise smoothness penalties to control overparameterized overfitting. 
Our proposed regularizers successfully mitigate Level 5 overfitting for both optimizers:
\begin{itemize}
    \item For the zero-order \textbf{1+1 ES} optimizer, enabling ESR and TV recovers Level 5 generalization performance from 29.43\% (unregularized) back to \textbf{30.17\%} (a substantial recovery of \textbf{0.74\%}), recovering nearly the entire drop compared to the unregularized L5 which collapses to 29.43\% and coming extremely close to the unoptimized uniform baseline of 30.41\%.
    \item For first-order \textbf{Adam}, the proposed regularizers recover Level 5 performance from 26.91\% to \textbf{28.51\%} (a massive recovery of \textbf{1.60\%}). While this is a clear and statistically robust stabilization effect, Adam remains more vulnerable to overfitting than ES, confirming that gradient descent in high dimensions requires strong explicit structural constraints (e.g., piece-wise polynomial or spline parameterizations) rather than simple soft penalties.
\end{itemize}

\textbf{4. The Supremacy of Static Baselines and Surrogate Loss Misalignment.} 
A key, high-signal finding of our empirical investigation is that despite optimization, no test-time adaptive merging configuration (even when regularized with ESR and TV) outperforms the static, zero-overhead Uniform Task Arithmetic baseline of 30.41\%. This phenomenon holds across both first-order and zero-order optimizers, and warrants a deep diagnostic analysis. At the core of this behavior lies a fundamental misalignment between the unsupervised surrogate loss (prediction entropy) and downstream classification accuracy under small calibration budgets ($N=256$). 

While prediction entropy serves as a convenient differentiable proxy for model confidence, minimizing it does not guarantee a correction of decision boundaries; instead, the optimizer often finds degenerate, "confident but incorrect" parameter configurations that yield sharp, low-entropy predictions for the incorrect classes. Under high parameter resolution (such as L5 Tensor-wise), this misalignment is severely amplified, as the optimizer has enough degrees of freedom to fit the transductive noise of the local calibration batch without learning generalizable features. For adaptive model merging to truly surpass static uniform blending in low-resource edge scenarios, future work must explore: (1) semantically richer surrogate objectives that incorporate contrastive or self-supervised losses rather than simple marginal entropy, (2) stronger structural priors (e.g., piece-wise polynomial or spline parameterizations) that restrict coefficient updates to low-frequency manifold subspaces, and (3) joint calibration over multi-modal distributions to prevent local transductive shift.

\subsection{Limitations and Scope}
To ensure absolute empirical transparency, we explicitly scope our experimental constraints:
\begin{itemize}
    \item \textbf{Resource-Constrained Scenario}. Due to strict CPU execution boundaries, we evaluate compact \texttt{ViTTiny} models ($d_{\text{model}}=64$) trained on relatively small task subsets (500 samples per task). While the absolute expert accuracies (61.03\% MNIST, 62.47\% FashionMNIST, 24.93\% CIFAR-10, 17.50\% SVHN) are lower than fully scaled pre-trained models on massive datasets, they represent a highly challenging, non-fully-converged \emph{low-resource edge warm-start setting} which serves as a rigorous, amplified "stress test" for deconstructing transductive overfitting dynamics.
    \item \textbf{Generalizability to High-Fidelity Experts}. We acknowledge that in poorly converged or low-fidelity experts, task vectors contain higher levels of high-frequency parameter noise. This noise potentially amplifies vulnerability to transductive overfitting during test-time adaptation. Whether fully converged, high-fidelity expert models (where features and representations are highly structured and stable) exhibit a similar sharp collapse at finer granularities, or whether high-fidelity representations are naturally more resilient to local transductive perturbations, remains an important open question. We strongly advocate for future work to validate these granularity-generalization dynamics on fully converged, high-fidelity models and large-scale foundations (e.g., CLIP-Large, LLaMA) deployed on massive test-time streams.
    \item \textbf{Evaluation Scale}. By evaluating on 1000 test samples per task (4000 samples overall), our test accuracy metrics achieve extremely high statistical robustness, ensuring that all macro-level differences and regularization recoveries are highly significant.
\end{itemize}
